\section{Introduction}\label{sec:intro}

Groups have to decide whom to listen to. A team follows whoever has the freshest information about a client; a
community of experts defers to whoever has been right lately; a flock follows the animals that seem to know where the
food is. In each case the source of useful information moves, and the signals about where it is are noisy. A group that
commits to a hub gains when the hub is right and pays when it is not; a group that never commits gains little but risks
little. This paper asks when such coordination pays, what it buys, what it costs, and what breaks it.

We study a minimal model. $N$ agents each split one unit of attention among the others. Every round a piece of news
appears at one agent, the source; with probability $\nu$ it appears at a random agent instead (a false alarm). Each agent
then strengthens the links that would have brought it the news soonest, counting paths through the network, with
self-reinforcement $\kappa$ (strong links earn more), commitment $T$ (how hard it bets on its best link) and memory
$\varepsilon$ (how fast old attention is replaced). The source stays put for $\tauE$ rounds on average and then moves.

The model joins three lines of work. The source-to-structure model of \citet{bromberg2020hierarchy} asked how the place
where information appears shapes who listens to whom, measuring closeness by a random-walk effective distance
\citep{brockmann2013hidden} and shape by hierarchy coordinates \citep{corominas2013originsheirarchy}. Optimal learning
in noisy niches \citep{lee2024constructing} asked how much memory a learner should keep when its environment switches,
and what happens when learners change the environment they track. Adaptive transport networks
\citep{corson2010fluctuations,katifori2010damage,ronellenfitsch2016global} grow the links that carry flow under a
budget, and form trees or loops depending on how the load fluctuates. No familiarity with any of them is assumed; the
model is defined from scratch in \cref{sec:model}, and \cref{sec:origins} says how each original premise fared.

\paragraph{Contributions.}
\begin{enumerate}[label=(\roman*)]
  \item A computed regime map with exact boundaries: lock-in for $T<\kappa$, coordinated hierarchy for
    $\kappa<T<\Thalf$, spread attention above $\Thalf$ (\cref{sec:regimes}). The one-variable reduction behind the
    boundaries is exact at any walk cost (A1), it remains the linear stability of the full network despite a kink at the
    fixed point (A2), and lock-in below $\kappa$ is decided by a separatrix in the initial weights (A3).
  \item Benefit and cost of coordination as functions of group size and noise (\cref{sec:scaling}): structural cost is
    flat in $N$ while benefit grows like $\ln(N-1)$; noise narrows the usable window of commitment instead of raising the
    cost.
  \item A mechanism for noise robustness near the lock-in edge, isolated by a factorial test (\cref{sec:robustness}):
    unequal rivals compete less because the score is concave in weight, and the hub relays even false news.
  \item A comparison of reward rules (\cref{sec:learning}): position (network closeness) against direct evidence
    (likelihood) and an ideal Bayesian observer.
  \item Results on memory, environmental pace and paid attention (\cref{sec:memory}), and on budgets and network shape
    against transport networks (\cref{sec:structure}).
\end{enumerate}

A companion website presents the same material as interactive chapters with enlargeable figures. All simulations,
analyses and checks are reproducible from the repository (\cref{app:methods}).
